\section{The small-factor transfer}
\label{sec:small-factor-transfer}

We next evaluate \eqref{eq:intro-small-residual}.  Here $n$ is the modulus
seen by the target prime.  This is a one-sided Type~I problem, so the
classical square-root level applies.

\begin{theorem}[Coefficient-uniform fixed-product Type~I estimate]
\label{thm:fixed-product-Type-I}
Fix $h\ne0$, $W\in C_c^1((0,\infty))$, and $A>0$.  There is
$B=B(A,h,W)>0$ such that, whenever
\begin{equation}\label{eq:fixed-product-BV-range}
 1\leq D\leq\frac{X^{1/2}}{(\log X)^B},
\end{equation}
one has, uniformly for arbitrary coefficients $|a_n|\leq1$,
\begin{equation}\label{eq:fixed-product-Type-I}
 \sum_{\substack{n\leq D\\(n,h)=1}}a_n
 \left\{
   \sum_{m\geq1}\Lambda(nm+h)W(nm/X)
   -\frac{X}{\varphi(n)}I_0(W)
 \right\}
 \ll_{A,h,W}\frac{X}{(\log X)^A}.
\end{equation}
\end{theorem}

\begin{proof}
For $(n,h)=1$, the substitution $t=nm+h$ places the target in the reduced
class $h\pmod n$ and changes the weight to $W((t-h)/X)$.  Stieltjes partial
summation bounds the row error by a constant depending on $W$ times
\[
 \max_{y\ll_W X+|h|}
 \left|
  \sum_{\substack{t\leq y\\t\equiv h\pmod n}}\Lambda(t)
  -\frac{y}{\varphi(n)}
 \right|.
\]
The continuous main term is exactly
\[
 \frac1{\varphi(n)}\int_0^\infty W((t-h)/X)\,dt
 =\frac{X}{\varphi(n)}I_0(W).
\]
Sum the maximal row errors over $n\leq D$ and apply the classical maximal
Bombieri--Vinogradov theorem with a slightly larger requested logarithmic
saving \citep{Bombieri1965,Vinogradov1965,IwaniecKowalski2004}.
\end{proof}

Rows not coprime to $h$ are supported on a finite collection of
prime-power targets.  We isolate this elementary point because the source
coefficient is now $\Lambda(n)-1$, rather than a M\"obius divisor weight.

\begin{lemma}[Noncoprime residual rows]
\label{lem:noncoprime-residual-rows}
For every $D\leq2X$ and every $\eps>0$,
\begin{equation}\label{eq:noncoprime-residual-rows}
 \sum_{\substack{n\leq D\\(n,h)>1}}
 |\Lambda(n)-1|
 \sum_{m\geq1}\Lambda(nm+h)|W(nm/X)|
 \ll_{h,W,\eps}X^\eps.
\end{equation}
Without an artificial upper cutoff, one also has
\begin{equation}\label{eq:noncoprime-residual-tail}
 \sum_{\substack{m,n\geq1\\n>D\\(n,h)>1}}
 |\Lambda(n)-1|\Lambda(nm+h)|W(nm/X)|
 \ll_{h,W,\eps}X^\eps.
\end{equation}
\end{lemma}

\begin{proof}
For either displayed sum, put $r=nm$ and enlarge the divisor range.  The
relevant left-hand side is at most
\[
 \sum_{r\asymp_W X}\Lambda(r+h)|W(r/X)|
 \sum_{\substack{n\mid r\\(n,h)>1}}(\Lambda(n)+1).
\]
If the inner sum is nonempty, some prime $p\mid h$ divides $r$.  Then
$p\mid r+h$.  If $\Lambda(r+h)\ne0$, it follows that $r+h=p^k$.
For each of the finitely many primes dividing $h$, only $O_{h,W}(1)$ powers
lie in the fixed-ratio interval forced by the support of $W$.  Finally,
\[
 \sum_{n\mid r}(\Lambda(n)+1)=\log r+\tau(r)\ll_\eps X^\eps,
\]
after reducing $\eps$ to absorb the remaining logarithms.
\end{proof}

\begin{theorem}[Small-factor prime-residual transfer]
\label{thm:small-factor-residual}
Fix $A>0$ and choose $B=B(A,h,W)$ sufficiently large.  Uniformly for
\begin{equation}\label{eq:small-factor-D-range}
 2\leq D\leq\frac{X^{1/2}}{(\log X)^B},
\end{equation}
one has
\begin{equation}\label{eq:small-factor-residual}
 \cA_{h,\leq D}(X;W)
 =XI_0(W)
  \sum_{\substack{n\leq D\\(n,h)=1}}
  \frac{\Lambda(n)-1}{\varphi(n)}
 +O_{A,h,W}\!\left(\frac{X}{(\log X)^A}\right).
\end{equation}
\end{theorem}

\begin{proof}
The noncoprime rows are absorbed by
\cref{lem:noncoprime-residual-rows}.  On the remaining rows apply
\cref{thm:fixed-product-Type-I} with
\[
 a_n=\frac{\Lambda(n)-1}{1+\log D}.
\]
These coefficients have modulus at most one.  Multiply the result back by
$1+\log D$ and request $A+3$ logarithmic saving in the Type~I theorem.
This proves \eqref{eq:small-factor-residual}.
\end{proof}

\begin{remark}[Why the coefficient is permitted here]
The source factor $\Lambda(n)-1$ is unbounded and nonperiodic, but in
\cref{thm:small-factor-residual} it weights the \emph{modulus} average.
After division by $1+\log D$, maximal Bombieri--Vinogradov is uniform in
that row coefficient.  This does not create a free second coefficient
inside an individual prime progression and therefore does not constitute a
Type~II estimate.
\end{remark}
